% =====================================================================================
% PREÁMBULO DE LA PLANTILLA (paquetes y estilos)  -->  NORMALMENTE NO NECESITAS EDITARLO
% Actualización: Said Polanco-Martagón / Marco A. Nuño-Maganda. Reorganizado para el sistema de revisión.
% =====================================================================================

\usepackage[left=2.5cm,top=2.0cm,right=2.5cm,bottom=3.0cm]{geometry}
\usepackage[utf8]{inputenc}
\usepackage[T1]{fontenc}
\usepackage[spanish,mexico]{babel}
\usepackage{etoolbox}
\usepackage{enumitem}                    % listas compactas en las guías
\usepackage[linguistics]{forest}
\usepackage{tikz}
\usepackage{amssymb, amsmath, amsbsy}   % símbolos y ecuaciones
\usepackage{longtable}                  % tablas largas
\usepackage{graphicx}
\usepackage{fancyhdr}
\usepackage{xcolor}
\usepackage{multirow}
\usepackage{listings}
\usepackage{float}                      % necesario para la opción [H] de figuras y tablas
\usepackage{caption}
\usepackage{subcaption}
\usepackage[skip=12pt plus1pt]{parskip}
\usepackage{pdfpages}                   % incluir PDF (cartas) en el documento
\usepackage{verbatim}
\usepackage{algpseudocode}
\usepackage{algorithm}
\usepackage{pdflscape}
\usepackage{afterpage}
\usepackage{array,booktabs,ragged2e}    % booktabs: \toprule \midrule \bottomrule en las tablas
\usepackage{wallpaper}
\usepackage[overload]{textcase}
\usepackage{datetime}

\newcolumntype{L}[1]{>{\raggedright\let\newline\\\arraybackslash\hspace{0pt}}m{#1}}
\newcolumntype{C}[1]{>{\centering\let\newline\\\arraybackslash\hspace{0pt}}m{#1}}
\newcolumntype{R}[1]{>{\raggedleft\let\newline\\\arraybackslash\hspace{0pt}}m{#1}}

\floatname{algorithm}{Algoritmo}
\renewcommand{\listalgorithmname}{Índice de algoritmos}
\renewcommand{\algorithmicrequire}{\textbf{Entrada:}}
\renewcommand{\algorithmicensure}{\textbf{Salida:}}

% ---------------------------------------------------------------------------------------
% REFERENCIAS (biblatex + biber). El estilo se elige en configuracion/datos.tex (\EstiloCitas).
% Archivo de bibliografía: referencias/referencias.bib
% ---------------------------------------------------------------------------------------
\usepackage{csquotes}
\ifdefstring{\EstiloCitas}{apa}
  {\usepackage[backend=biber,style=apa,sorting=nyt]{biblatex}
   \DeclareLanguageMapping{spanish}{spanish-apa}
   % Ajustes adicionales SOLO para el estilo de citas APA (se carga si \EstiloCitas = apa).
% Mejoran el uso de «y» en español y el separador de nombres en la lista de referencias.
\makeatletter
\DefineBibliographyExtras{spanish}{%
  \setcounter{smartand}{1}%
  \let\lbx@finalnamedelim=\lbx@es@smartand
  \let\lbx@finallistdelim=\lbx@es@smartand
}
\makeatother
}
  {\usepackage[backend=biber,style=ieee,sorting=none]{biblatex}}
\addbibresource{referencias/referencias.bib}
\DefineBibliographyStrings{spanish}{references = {Referencias}, bibliography = {Referencias}}

% Permite romper los hipervínculos largos al final de la línea
\setcounter{biburllcpenalty}{7000}
\setcounter{biburlucpenalty}{8000}

\usepackage[bookmarks=true,breaklinks=true,bookmarksopen=false,colorlinks=true,linkcolor=blue]{hyperref}

% Fecha de compilación (se usa en las citas de páginas de Internet)
\newdateformat{specialdate}{\twodigit{\THEDAY}-\twodigit{\THEMONTH}-\THEYEAR}
\date{\specialdate\today}

\newcommand{\HRule}{\rule{\linewidth}{0.25mm}}
\newcommand{\separacionCorta}{0.0cm}
\newcommand{\separacionLarga}{0.5cm}
\newcommand{\iemph}[1]{\MakeTextUppercase{#1}}

% ---------------------------------------------------------------------------------------
% ENCABEZADO Y PIE DE PÁGINA
% ---------------------------------------------------------------------------------------
\pagestyle{fancy}
\setlength{\headheight}{50pt}
\fancyhead{}
\fancyhead[L]{\includegraphics[height=1.00cm]{logos/UTYP.png}}
\fancyhead[C]{\begin{center}\scriptsize{\NombreProyectoheader}\end{center}}
\fancyhead[R]{\includegraphics[height=1.25cm]{logos/LogoUPV_2023.png}}
\fancyfoot{}
\fancyfoot[R]{\thepage}

% ---------------------------------------------------------------------------------------
% GUÍAS PARA EL ALUMNO
% Las guías (entorno guia) explican qué debe ir en cada sección y qué preguntas debe responder.
%   * \guiastrue   -> se muestran en un recuadro (valor inicial, mientras escribes).
%   * \guiasfalse  -> se ocultan en el PDF (cámbialo para la versión final).
% El sistema de revisión IGNORA todo lo que esté dentro de \begin{guia} ... \end{guia}: lo que escribas ahí no
% cuenta como contenido de tu tesina. Escribe tu texto FUERA de las guías.
% ---------------------------------------------------------------------------------------
\usepackage[most]{tcolorbox}
\newif\ifguias
\guiastrue
\newtcolorbox{guiabox}[1][]{breakable, enhanced jigsaw, colback=yellow!10, colframe=orange!70!black,
  fonttitle=\bfseries\small, title={#1}, boxrule=0.6pt, arc=2pt, left=5pt, right=5pt, top=3pt, bottom=3pt,
  fontupper=\small, before skip=8pt, after skip=10pt}
\newenvironment{guia}[1][Guía para el alumno]
  {\ifguias\begin{guiabox}[#1]\else\setbox0\vbox\bgroup\fi}
  {\ifguias\end{guiabox}\else\egroup\fi}
